% Hafta 13 -- alıştırma gövdesi (ders notu ve ayrı alıştırma kağıdı tarafından \input edilir)
\section{Çözümlü Alıştırmalar}
Her soruyu önce kendiniz deneyin, sonra çözümü okuyun.

\begin{alistirma}
Ardışık üç çift sayının toplamı \(138\)'dir. En küçüğü kaçtır?
\end{alistirma}
\begin{cozum}
Sayılar \(x,\ x+2,\ x+4\) (\(x\) çift):
\[x+(x+2)+(x+4)=138\ \Rightarrow\ 3x+6=138\ \Rightarrow\ x=44\]
Sayılar \(44,\ 46,\ 48\); toplam \(138\) \checkmark. \ (Kısa yol: ortanca \(138\div 3=46\).)
\end{cozum}

\begin{alistirma}
Toplamları \(85\) olan iki doğal sayıdan büyüğü küçüğüne bölündüğünde bölüm \(3\), kalan \(9\) oluyor. Büyük sayı kaçtır?
\end{alistirma}
\begin{cozum}
Küçük sayı \(k\) ise büyük sayı \(3k+9\) (bölme: bölünen \(=\) bölen \(\cdot\) bölüm \(+\) kalan).
\[k+3k+9=85\ \Rightarrow\ 4k=76\ \Rightarrow\ k=19,\qquad \text{büyük}=3\cdot 19+9=66\]
Sağlama: \(66=19\cdot 3+9\) ve kalan bölenden küçük (\(9<19\)) \checkmark
\end{cozum}

\begin{alistirma}
İki basamaklı bir sayının rakamları toplamı \(9\)'dur. Rakamlarının yerleri değiştirilince sayı \(45\) azalıyor. Bu sayı kaçtır?
\end{alistirma}
\begin{cozum}
\(\ab{ab}-\ab{ba}=(10a+b)-(10b+a)=9(a-b)=45\Rightarrow a-b=5\).
\[a+b=9,\ \ a-b=5\ \Rightarrow\ 2a=14\ \Rightarrow\ a=7,\ b=2\ \Rightarrow\ \ab{ab}=72\]
Sağlama: \(72-27=45\) \checkmark
\end{cozum}

\begin{alistirma}
Bir kişi parasının \(\frac{3}{8}\)'ini harcıyor, kalan parasının \(\frac{2}{5}\)'ini de kardeşine veriyor. Geriye \(180\) TL kaldığına göre başlangıçta kaç TL'si vardı?
\end{alistirma}
\begin{cozum}
Paydalar \(8\) ve \(5\) \(\Rightarrow\) para \(=40k\).\\
Harcanan: \(\frac{3}{8}\cdot 40k=15k\) \(\to\) kalan \(25k\); \ kardeşine: \(\frac{2}{5}\cdot 25k=10k\) \(\to\) kalan \(15k\).
\[15k=180\ \Rightarrow\ k=12\ \Rightarrow\ 40k=480\ \text{TL}\]
Sağlama: \(480-180=300\), \ \(300-120=180\) \checkmark
\end{cozum}

\begin{alistirma}
Bir su deposunun \(\frac{1}{3}\)'ü doludur. Depoya \(50\) litre su eklenince deponun \(\frac{3}{4}\)'ü doluyor. Deponun tamamı kaç litre su alır?
\end{alistirma}
\begin{cozum}
Eklenen su, doluluk oranlarının farkına karşılık gelir:
\[\frac{3}{4}-\frac{1}{3}=\frac{9-4}{12}=\frac{5}{12}\ \Rightarrow\ \frac{5}{12}V=50\ \Rightarrow\ V=50\cdot\frac{12}{5}=120\ \text{L}\]
Sağlama: \(\frac{1}{3}\cdot 120=40\), \ \(40+50=90=\frac{3}{4}\cdot 120\) \checkmark
\end{cozum}

\begin{alistirma}
İki kardeşin yaşları toplamı \(30\)'dur. \(5\) yıl önce büyük kardeşin yaşı küçüğün yaşının \(3\) katıydı. Kardeşler bugün kaç yaşındadır?
\end{alistirma}
\begin{cozum}
Tablo yöntemi (küçüğün yaşı \(k\), büyüğünki \(30-k\)):
\begin{center}
\begin{tabular}{lcc}
 & \(5\) yıl önce & Şimdi\\\hline
Küçük & \(k-5\) & \(k\)\\
Büyük & \(25-k\) & \(30-k\)
\end{tabular}
\end{center}
\[25-k=3(k-5)\ \Rightarrow\ 25-k=3k-15\ \Rightarrow\ 4k=40\ \Rightarrow\ k=10\]
Küçük \(10\), büyük \(20\) yaşında. \quad Sağlama: \(5\) yıl önce \(15=3\cdot 5\) \checkmark
\end{cozum}

\begin{alistirma}
Bir baba \(46\) yaşında, üç çocuğu \(14\), \(10\) ve \(6\) yaşındadır. Kaç yıl sonra babanın yaşı çocuklarının yaşları toplamına eşit olur?
\end{alistirma}
\begin{cozum}
Çocukların yaş toplamı bugün \(30\); her yıl \(3\) artar (\(3\) kişi). Babanın yaşı her yıl \(1\) artar.
\[46+t=30+3t\ \Rightarrow\ 2t=16\ \Rightarrow\ t=8\]
Sağlama: \(8\) yıl sonra baba \(54\); çocuklar \(22+18+14=54\) \checkmark
\end{cozum}

\begin{alistirma}
Mert \(32\) yaşındadır. Mert, Ozan'ın bugünkü yaşındayken Ozan \(20\) yaşındaydı. Ozan, Mert'in bugünkü yaşına geldiğinde Mert kaç yaşında olur?
\end{alistirma}
\begin{cozum}
Ozan \(x\), yaş farkı \(d=32-x\). Mert'in \(x\) yaşında olduğu an \(d\) yıl öncedir; o zaman Ozan \(x-d=20\):
\[x-(32-x)=20\ \Rightarrow\ 2x=52\ \Rightarrow\ x=26,\quad d=6\]
Ozan'ın \(32\) yaşına gelmesi \(6\) yıl sürer; o zaman Mert \(32+6=38\) yaşında olur.\\
Sağlama: \(20,\ 26,\ 32,\ 38\) eşit aralıklı (\(d=6\)) \checkmark
\end{cozum}

\kendinizoncesi
\section{Alıştırmalar (Kendiniz Çözün)}

\begin{enumerate}[label=\textbf{\arabic*.},leftmargin=*,itemsep=6pt]
\item Bir sayının \(4\) katının \(7\) fazlası \(43\)'tür. Bu sayı kaçtır?
\item Bir sayının \(3\) eksiğinin \(5\) katı, aynı sayının \(2\) katının \(6\) fazlasına eşittir. Bu sayı kaçtır?
\item İki sayının toplamı \(75\), farkı \(21\)'dir. Küçük sayı kaçtır?
\item Ardışık beş tek sayının toplamı \(115\)'tir. En büyüğü kaçtır?
\item \(\ab{ab}\) ve \(\ab{ba}\) iki basamaklı sayılardır. \(\ab{ab}+\ab{ba}=154\) ve \(a-b=2\) ise \(\ab{ab}\) kaçtır?
\item Bir sınıftaki öğrenciler sıralara \(4\)'erli oturursa \(2\) sıra boş kalıyor (diğer sıralar tam dolu); \(3\)'erli oturursa \(5\) öğrenci ayakta kalıyor. Sınıf mevcudu kaçtır?
\item Bir sayının \(\frac{3}{5}\)'i \(27\) ise bu sayının \(\frac{2}{9}\)'u kaçtır?
\item Bir kişi parasının \(\frac{1}{6}\)'sını harcıyor, kalan parasının \(\frac{2}{5}\)'ini de arkadaşına veriyor. Geriye \(90\) TL kaldığına göre başlangıçta kaç TL'si vardı?
\item Bir deponun \(\frac{5}{8}\)'i doludur. Depodan \(36\) litre su kullanılınca deponun \(\frac{1}{4}\)'ü dolu kalıyor. Deponun tamamı kaç litredir?
\item Bir sayının \(\frac{3}{4}\)'ünün \(2\) eksiği, aynı sayının \(\frac{2}{3}\)'ünün \(1\) fazlasına eşittir. Bu sayı kaçtır?
\item Bir anne \(34\), oğlu \(7\) yaşındadır. Kaç yıl sonra annenin yaşı oğlunun yaşının \(4\) katı olur?
\item \(6\) kişilik bir ailenin yaş ortalaması \(20\)'dir. Ailenin en küçük üyesi \(7\) yaşında olduğuna göre \(5\) yıl önce bu ailenin yaşları toplamı kaçtı?
\item Bir babanın yaşı kızının yaşının \(5\) katıdır. \(3\) yıl sonra babanın yaşı kızının yaşının \(4\) katı olacaktır. Baba bugün kaç yaşındadır?
\item Can, Ece'den \(6\) yaş büyüktür. Ece, Can'ın bugünkü yaşına geldiğinde ikisinin yaşları toplamı \(50\) olacaktır. Ece bugün kaç yaşındadır?
\item Bir baba 1992'de, oğlu 2019'da doğmuştur. Hangi yılda babanın yaşı oğlunun yaşının \(4\) katı olur?
\end{enumerate}

\subsection*{Cevap Anahtarı}
{\small
\begin{multicols}{2}
\begin{enumerate}[label=\textbf{\arabic*.},leftmargin=*,itemsep=1pt]
\item \(4x+7=43\Rightarrow x=9\).
\item \(5x-15=2x+6\Rightarrow x=7\).
\item \((75-21)/2=27\) \ (büyük \(48\)).
\item Ortanca \(115/5=23\); sayılar \(19,\dots,27\) \(\Rightarrow\) \(27\).
\item \(11(a+b)=154\Rightarrow a+b=14\); \(a=8,b=6\) \(\Rightarrow\) \(86\).
\item \(4(x-2)=3x+5\Rightarrow x=13\) sıra; mevcut \(44\).
\item Sayı \(27\cdot\frac{5}{3}=45\); \(\frac{2}{9}\cdot 45=10\).
\item Para \(30k\): \(5k\), kalan \(25k\); \(10k\), kalan \(15k=90\Rightarrow 180\) TL.
\item \(\frac{5}{8}-\frac{1}{4}=\frac{3}{8}\); \(\frac{3}{8}V=36\Rightarrow V=96\) L.
\item \(\cdot 12\): \(9x-24=8x+12\Rightarrow x=36\).
\item \(34+t=4(7+t)\Rightarrow t=2\) \ (\(36=4\cdot 9\)).
\item Bugün \(120\); \(120-6\cdot 5=90\).
\item \(5x+3=4(x+3)\Rightarrow x=9\); baba \(45\).
\item \(6\) yıl sonra: Ece \(E+6\), Can \(E+12\); \(2E+18=50\Rightarrow E=16\).
\item Fark \(27\); \(4k=k+27\Rightarrow k=9\Rightarrow 2028\).
\end{enumerate}
\end{multicols}}
